% =====================================================================
%  chapter01_psychology.tex — הפסיכולוגיה של הכלב העירוני
% =====================================================================
\chapter{הפסיכולוגיה של הכלב העירוני}
\chaptersubtitle{איך לשרוד מעלית עם שלושה כלבים וחתול}

\section{העיר מגובה חצי מטר}
נסו לרדת לרגע לגובה של כלב בינוני. המכוניות נראות פתאום כמו קירות נעים, הצופר של האוטובוס
מגיע בדיוק בגובה האוזניים, ורגליים של זרים חולפות במרחק נגיעה. מעליכם מתנהלת שיחה שאתם לא
מבינים, ומתחתיכם מדרכה שמספרת סיפור ריחני על כל מי שעבר בה היום.

לחלק מהכלבים זה גן עדן. אחרים פשוט צריכים יותר זמן להתרגל, ויש כלבים שהעומס הזה מתיש אותם.
אין כאן "נכון" אחד: גנטיקה, ניסיון מוקדם, גיל ובריאות משפיעים על האופן שבו כל כלב חווה את
הרחוב. המשימה שלנו היא ללמוד לקרוא את הכלב \emph{שלנו}, לא את הכלב הממוצע.

\section{רשימת הגירויים העירונית}
מכוניות וצפירות, אופניים וקורקינטים חשמליים שמופיעים בלי רעש, מעליות שנסגרות בחבטה, ילדים
שרצים לחבק, כלבים אחרים מעבר לכביש ושקית פלאפל שמישהו שכח על הספסל. כל אחד מהגירויים האלה
לבדו הוא עניין קטן. כשהם מצטברים בטיול אחד, חלק מהכלבים מגיעים לעומס חושי: הם מתקשים להקשיב,
מושכים, נובחים או "נתקעים" במקום.

\begin{knowbox}
התנהגות "קשה" ברחוב היא לפעמים לא חוסר חינוך אלא סימן לכך שהכלב מוצף. ארגוני רווחת בעלי חיים
ממליצים ללמוד את שפת הגוף של הכלב כדי לזהות לחץ מוקדם, לפני שהוא מסלים~\cite{rspca-behaviour}.
\end{knowbox}

\section{שפת הגוף: מה הכלב מנסה להגיד}
כלבים "מדברים" בעיקר בגוף. הטבלה הקטנה שבהמשך היא נקודת פתיחה בלבד: תמיד מסתכלים על \emph{כל}
הגוף ועל ההקשר, ולא על סימן בודד. גם זנב מכשכש לא אומר בהכרח שהכלב שמח: לפעמים כך נראית
דווקא התרגשות מתוחה.

\begin{tcbraster}[raster columns=3,raster equal height,raster column skip=6pt,
   raster before skip=6pt,raster after skip=8pt,
   raster every box/.append style={before upper={\raggedright}}] % ragged end side (DET-1-08)
  \begin{tcolorbox}[enhanced,colback=sage!14!white,colframe=sagedark,arc=2mm,boxrule=0.8pt,
     title={\sffamily\bfseries רגוע ונינוח},colbacktitle=sagedark,fontupper=\small,left=4pt,right=4pt]
    גוף רפוי, פה פתוח מעט, תנועות "גליות", זנב בגובה טבעי, מבט רך.
  \end{tcolorbox}
  \begin{tcolorbox}[enhanced,colback=gold!16!white,colframe=golddark,arc=2mm,boxrule=0.8pt,
     title={\sffamily\bfseries סקרן או נרגש},colbacktitle=golddark,fontupper=\small,left=4pt,right=4pt]
    אוזניים קדימה, גוף נוטה לעבר הגירוי, ריחרוח ממוקד. התרגשות יכולה להיראות כמו קפיצות ומשיכה.
  \end{tcolorbox}
  \begin{tcolorbox}[enhanced,colback=coral!14!white,colframe=coraldark,arc=2mm,boxrule=0.8pt,
     title={\sffamily\bfseries לחוץ או מפחד},colbacktitle=coraldark,fontupper=\small,left=4pt,right=4pt]
    זנב מקופל, אוזניים לאחור, ליקוק שפתיים, פיהוק "לא במקום", לובן עין גלוי, קפיאה או ניסיון לברוח~\cite{rspca-behaviour,aspcapro-body}.
  \end{tcolorbox}
\end{tcbraster}

ההבדל בין התרגשות לפחד לא תמיד ברור, וזה בסדר. אם אתם לא בטוחים, תנו לכלב יותר מרחק מהגירוי
ובדקו אם הגוף שלו משתחרר. כלב שנרגע כשמתרחקים כנראה היה בלחץ.

\section{חיזוק חיובי ושגרה צפויה}
הדרך היעילה ביותר לעזור לכלב עם העיר היא ללמד אותו שדברים מפחידים מנבאים דברים טובים: חטיף
כשהאוטובוס עובר, שבח כשהוא מסתכל עליכם במקום על הכלב שממול. איגוד הווטרינרים האמריקאי לרפואה
התנהגותית (\en{AVSAB}) ממליץ על אימון מבוסס תגמול ומתנגד לשימוש בשיטות ענישה~\cite{avsab2021}, וסקירת
מחקר וסקר בקרב בעלי כלבים מצאו ששיטות אימון מבוססות ענישה או הפחדה קשורות לסיכון גבוה יותר
לבעיות התנהגות ולפגיעה ברווחה~\cite{ziv2017,hiby2004}.

שגרה עוזרת לא פחות. כלב שיודע מתי יוצאים, באיזה מסלול ומה קורה בסוף, מתמודד בקלות רבה יותר
עם ההפתעות שהעיר תמיד מספקת.

\section{מעלית: שלושה כלבים, חתול אחד וארבע קומות של מתח}

\begin{figure}[h]
  \centering
  \elevatorScene
  \caption{"מעלית לקומה 7". איור מקורי שנבנה בקוד \en{TikZ} מצורות גאומטריות (כלב, חתול, מעלית).}
  \label{fig:elevator}
\end{figure}

קומה 3. הדלת נפתחת. בפנים: רקסי, בורדר קולי עם תוכניות; פיצי, צ\gr יוואווה עם דעות; ומיסטר ביג,
דני ענק שמנסה להיראות קטן. בפינה יושב החתול של הקומה השביעית, שרכב במעלית הזאת עוד לפני
שמישהו מהנוכחים נולד. כולם מסתכלים על כולם. אף אחד לא נושם.

הנה איך זה \emph{לא} נגמר בבלגן:

\begin{tipbox}[כללי המעלית]
\begin{itemize}
  \item \textbf{שואלים לפני שנכנסים.} "זה בסדר אם אצטרף עם הכלב?" זאת לא חולשה, זה נימוס.
  \item \textbf{רצועה קצרה, לא מתוחה.} הכלב צמוד אליכם, והרצועה רפויה ולא נכרכת סביב רגליים.
  \item \textbf{פנים אל הדלת, גב אל הקיר.} מקטינים מבטים ישירים בין בעלי החיים.
  \item \textbf{מחכים למעלית הבאה} כשהכלב שלכם, או של מישהו אחר, נראה לחוץ. ההמתנה קצרה, התקרית
        ארוכה.
  \item \textbf{מתגמלים שקט.} חטיף קטן על ישיבה רגועה הופך את המעלית ממלכודת לחוויה טובה.
\end{itemize}
\end{tipbox}

\section{מפגשים עם כלבים אחרים ומרחב אישי}
לא כל כלב רוצה להכיר כל כלב, בדיוק כמו שלא כל אדם רוצה לדבר עם כל מי שעומד איתו בתור. מפגש
"אף מול אף" על מדרכה צרה, ברצועה, הוא מצב רגיש במיוחד: לאף אחד מהכלבים אין לאן לזוז.
שאלו את הבעלים לפני שמתקרבים, אפשרו לכלבים להיפגש בקשת ולא חזיתית, וסיימו את המפגש אחרי כמה
שניות, כשכולם עוד נהנים.

\begin{mistakebox}
"הוא ידידותי!" זה משפט נחמד, אבל הוא מתאר רק את הכלב \emph{שלכם}. ייתכן שהכלב השני חולה, מבוגר,
בשיקום או פשוט לא במצב רוח. מרחב אישי הוא זכות גם בעולם של ארבע רגליים.
\end{mistakebox}

\begin{fieldbox}
"שיטת הקשת": כשכלב מתקרב מולכם ואתם לא בטוחים, חצו לצד השני או עשו עיקוף קטן. הכלב שלכם
יקבל מרחק, ואתם תקבלו עוד שנייה לחשוב.
\end{fieldbox}

\section{גור בעיר: לאט, ובקצב של הכלב}
גורים לומדים על העולם במהירות, והחוויות של החודשים הראשונים משפיעות על האופן שבו יתמודדו
עם העיר בהמשך. איגוד הווטרינרים האמריקאי לרפואה התנהגותית ממליץ על חשיפה מוקדמת, מבוקרת
וחיובית לסביבות, לאנשים ולגירויים, בתיאום עם הווטרינר לגבי מצב החיסונים~\cite{avsab-puppy}.
המילה החשובה כאן היא \emph{הדרגתית:} לא "להטביע" את הגור ברעש של מרכז העיר, אלא לבנות סולם.

\begin{figure}[h]
  \centering
  \resizebox{0.97\linewidth}{!}{% =====================================================================
%  exposure_ladder.tex — "סולם החשיפה": gradual exposure to city stimuli (Chapter 1)
%  Five rising steps drawn from rounded rectangles; progression reads right -> left (Hebrew).
%  RTL: wrapped node text uses the `rtlpar` style, single-line labels use \RLE{} (see preamble).
% =====================================================================
\begin{tikzpicture}[
    font=\sffamily\small,
    ladderstep/.style={rtlpar,draw=none,rounded corners=2.5mm,align=center,text width=2.55cm,
                 inner sep=4pt,anchor=south,text=white},
    num/.style={circle,fill=white,text=navy,font=\sffamily\scriptsize\bfseries,inner sep=1.2pt,minimum size=13pt}
  ]
  % ground line
  \draw[sage!60!white,line width=1.2pt] (-0.4,0) -- (15.6,0);
  % steps: \i = position, colour mixes sage -> coral as difficulty rises
  \foreach \i/\head/\body in {
      1/צפייה מרחוק/מספסל ברחוב שקט,
      2/רחוב שקט/הליכה קצרה ומעט אנשים,
      3/רעשים מוכרים/אוטובוס עובר במרחק,
      4/מקום הומה/כמה דקות ליד בית קפה,
      5/מעלית ושכנים/נסיעה קצרה וחטיף בסוף}{
    \pgfmathsetmacro{\x}{14.0-3.15*(\i-1)}
    \pgfmathsetmacro{\h}{1.15+0.42*(\i-1)}
    \pgfmathsetmacro{\mix}{100-25*(\i-1)}
    \fill[sagedark!\mix!coraldark,rounded corners=2.5mm] (\x-1.42,0) rectangle (\x+1.42,\h);
    \node[ladderstep] at (\x,0.08) {\textbf{\head}\\[1pt]{\footnotesize \body}};
    \node[num] at (\x+1.3,\h) {\i};
  }
  % progression arrow (right -> left) above the steps
  \draw[-{Stealth[length=7pt]},navy,line width=1pt] (15.3,3.35) -- (-0.2,3.35);
  \node[fill=paper,inner sep=2pt,font=\sffamily\footnotesize\bfseries,text=navy] at (7.6,3.35)
       {\RLE{מתקדמים שלב רק כשהכלב רגוע בשלב הקודם}};
\end{tikzpicture}
}
  \caption{סולם החשיפה: דוגמה לחשיפה הדרגתית לגירויים עירוניים. כל שלב נמשך כמה ימים או שבועות,
  לפי תגובת הגור. איור מקורי בקוד \en{TikZ}.}
  \label{fig:ladder}
\end{figure}

בכל שלב מחפשים את אותם סימני נינוחות מהטבלה שלמעלה. אם הגור קופא, מנסה לברוח או מסרב לחטיף
שהוא בדרך כלל אוהב, זה סימן לחזור שלב אחד אחורה ולהגדיל את המרחק. הצלחה נמדדת בגור רגוע, לא
במספר הגירויים שהספקתם לסמן ביום אחד.

\section{לבד בבית: הדירה כשאתם בעבודה}
כלבים עירוניים רבים מבלים חלק מהיום לבד. חלקם ישנים בשלווה עד שאתם חוזרים, ואחרים מתקשים: נביחות
ממושכות, הרס ליד הדלת, ריור או צרכים בבית. חשוב לזכור שזו לא "נקמה": התנהגויות כאלה הן לרוב
סימן למצוקה כשהכלב נשאר לבד, ואפשר למנוע אותן ולטפל בהן~\cite{rspca-separation}.

\begin{tipbox}[לפני שיוצאים מהבית]
\begin{itemize}
  \item \textbf{טיול ממוקד לפני היציאה}, עם הרבה זמן ריחרוח, ולא רק "סיבוב להוריד".
  \item \textbf{העסקה:} צעצוע שאפשר למלא באוכל, משחק חיפוש או לעיסה מתאימה לגודל הכלב.
  \item \textbf{פרידה וחזרה שגרתיות:} בלי טקס ארוך ובלי התרגשות, כדי שהיציאה תהיה עניין רגיל.
  \item \textbf{תרגול הדרגתי:} מתחילים מהיעדרויות של דקות ספורות ומאריכים לאט.
\end{itemize}
\end{tipbox}

\begin{knowbox}[מתי לבקש עזרה]
אם השכנים מדווחים על נביחות ממושכות, או שהכלב פוגע בעצמו או בדלת, אל תחכו שזה "יעבור". פנו
לווטרינר כדי לשלול סיבה רפואית, ולאחר מכן למאלף או ליועץ התנהגות שעובד בשיטות מבוססות תגמול
\cite{avsab2021,rspca-separation}.
\end{knowbox}
